\documentclass[11pt,a4paper]{article}
\usepackage[utf8]{inputenc}
\usepackage[T1]{fontenc}
\usepackage{geometry}
\geometry{margin=1in}
\usepackage{hyperref}
\usepackage{enumitem}
\usepackage{booktabs}
\usepackage{longtable}
\usepackage{xcolor}
\usepackage{titlesec}
\usepackage{amsmath,amssymb}

\titleformat{\section}{\Large\bfseries\color{blue!70!black}}{}{0em}{}
\titleformat{\subsection}{\large\bfseries\color{blue!50!black}}{}{0em}{}
\titleformat{\subsubsection}{\normalsize\bfseries\color{blue!30!black}}{}{0em}{}

\title{\textbf{Replication Plan}\\[0.5em]
\large Approximating Photo-$z$ PDFs for Large Surveys\\[0.3em]
\normalsize OSTI ID: 1461824}
\author{Auto-generated Replication Blueprint}
\date{\today}

\begin{document}
\maketitle
\tableofcontents
\newpage

%---------------------------------------------------------------------
\section{Paper Overview}
%---------------------------------------------------------------------
This paper introduces \texttt{qp}, a Python package for efficiently approximating, storing, and manipulating photometric redshift probability density functions (photo-$z$ PDFs) for large astronomical surveys. The work evaluates multiple PDF approximation schemes---histograms, quantiles, samples, step functions, piecewise-linear interpolation, and Gaussian mixtures---using information-theoretic metrics (Kullback--Leibler divergence) and moments. The goal is to enable compact storage of billions of per-galaxy PDFs for surveys like LSST while preserving statistical fidelity.

%---------------------------------------------------------------------
\section{Detailed Workflow}
%---------------------------------------------------------------------

\subsection{Phase 1: Environment and Package Setup}
\begin{enumerate}[label=\textbf{1.\arabic*}]
  \item Clone the \texttt{qp} repository from GitHub.
  \item Install dependencies: NumPy, SciPy, scikit-learn, matplotlib, pathos (for parallelism).
  \item Verify installation by running the provided test suite.
\end{enumerate}

\subsection{Phase 2: Mock Photo-$z$ PDF Generation}
\begin{enumerate}[label=\textbf{2.\arabic*}]
  \item \textbf{Generate a library of mock photo-$z$ PDFs.}
    \begin{itemize}
      \item Use Gaussian mixture models with varying numbers of components (1--5) to simulate realistic multi-modal photo-$z$ PDFs.
      \item Sample redshift range $z \in [0, 3]$ with resolution $\Delta z = 0.01$.
      \item Generate $\sim$1000 diverse PDFs with varying widths, skewness, and multi-modality.
    \end{itemize}
  \item \textbf{Optionally use real photo-$z$ PDFs} from public catalogs (e.g., SDSS, BPZ outputs).
\end{enumerate}

\subsection{Phase 3: PDF Approximation Methods}
\begin{enumerate}[label=\textbf{3.\arabic*}]
  \item \textbf{Implement/use each approximation scheme:}
    \begin{itemize}
      \item \emph{Histogram}: bin the PDF into $N$ bins of uniform or variable width.
      \item \emph{Quantile}: store $N$ quantile values that summarize the CDF.
      \item \emph{Samples}: draw $N$ samples from the PDF.
      \item \emph{Step function}: piecewise-constant approximation.
      \item \emph{Piecewise linear}: linear interpolation between $N$ nodes.
      \item \emph{Gaussian mixture}: fit a $K$-component GMM to the PDF.
    \end{itemize}
  \item \textbf{Vary the compression parameter $N$} (number of bins/quantiles/samples/components) across a wide range.
\end{enumerate}

\subsection{Phase 4: Metric Computation}
\begin{enumerate}[label=\textbf{4.\arabic*}]
  \item \textbf{Compute KL divergence} $D_{\mathrm{KL}}(p \| \hat{p})$ between original and approximated PDFs.
  \item \textbf{Compute moment residuals}: mean, variance, skewness, kurtosis of original vs.\ approximated.
  \item \textbf{Compute stacked (ensemble) $n(z)$ residuals}: aggregate PDFs and compare the stacked distribution.
  \item \textbf{Record storage cost} (bytes per PDF) for each scheme at each $N$.
\end{enumerate}

\subsection{Phase 5: Results Reproduction}
\begin{enumerate}[label=\textbf{5.\arabic*}]
  \item \textbf{Reproduce metric-vs-$N$ curves} (Figures in the paper).
  \item \textbf{Reproduce storage-vs-fidelity trade-off plots.}
  \item \textbf{Reproduce the comparison table} of approximation schemes.
  \item \textbf{Validate} by comparing to paper's stated numerical results.
\end{enumerate}

%---------------------------------------------------------------------
\section{Datasets}
%---------------------------------------------------------------------

\begin{longtable}{p{4.5cm} p{3.5cm} p{6cm}}
\toprule
\textbf{Dataset / Source} & \textbf{Type} & \textbf{Location / Notes} \\
\midrule
\endfirsthead
\toprule
\textbf{Dataset / Source} & \textbf{Type} & \textbf{Location / Notes} \\
\midrule
\endhead
Mock photo-$z$ PDFs (Gaussian mixtures) & Synthetic & Generated using \texttt{qp}'s built-in utilities or custom scripts \\
\midrule
\texttt{qp} example/test data & Bundled test data & Included in the \texttt{qp} GitHub repository \newline \url{https://github.com/aimalz/qp} \\
\midrule
LSST DESC photo-$z$ challenge data (optional) & Public catalog & \url{https://github.com/LSSTDESC/RAIL} \\
\midrule
SDSS photometric redshift catalog (optional) & Public survey & \url{https://www.sdss.org/dr17/} \\
\midrule
BPZ photo-$z$ outputs (optional) & Public & Available via SDSS or DES data releases \\
\bottomrule
\end{longtable}

%---------------------------------------------------------------------
\section{Models, Tools, and Software}
%---------------------------------------------------------------------

\subsection{Programming Languages}
\begin{itemize}
  \item \textbf{Python 3.7+} --- all code is in Python.
\end{itemize}

\subsection{Libraries and Frameworks}
\begin{longtable}{p{4.5cm} p{2.5cm} p{6.5cm}}
\toprule
\textbf{Tool / Library} & \textbf{Version} & \textbf{Purpose / URL} \\
\midrule
\endfirsthead
\toprule
\textbf{Tool / Library} & \textbf{Version} & \textbf{Purpose / URL} \\
\midrule
\endhead
\texttt{qp} & latest & Core package for photo-$z$ PDF approximation \newline \url{https://github.com/aimalz/qp} \\
\midrule
NumPy & $\geq 1.18$ & Numerical array operations \\
\midrule
SciPy & $\geq 1.5$ & Statistical distributions, integration, interpolation \\
\midrule
scikit-learn & $\geq 0.24$ & Gaussian mixture model fitting \\
\midrule
Matplotlib & $\geq 3.3$ & Plotting metric curves and PDF comparisons \\
\midrule
pathos & $\geq 0.2.8$ & Parallel computation for large PDF ensembles \\
\midrule
Jupyter & $\geq 6.0$ & Notebook environment for interactive analysis \\
\bottomrule
\end{longtable}

\subsection{Hardware Requirements}
\begin{itemize}
  \item Any modern laptop/workstation; computations are lightweight.
  \item For scaling tests with millions of PDFs: $\geq$ 16\,GB RAM.
\end{itemize}

\end{document}
